\begin{itemize}
  \item VIV-32 defines sixteen 32-bit general-purpose registers, named \(r0\) through \(r15\).
  \item The ABI assigns conventional roles to these registers.
\end{itemize}

\begin{center}
\begin{tabular}{lll}
\hline
Register & ABI name & Role \\
\hline
r0  & zero & Hardwired zero register \\
r1  & rv0  & Return value 0 / scratch register \\
r2  & rv1  & Return value 1 / scratch register \\
r3  & a0   & Argument register 0 \\
r4  & a1   & Argument register 1 \\
r5  & a2   & Argument register 2 \\
r6  & a3   & Argument register 3 \\
r7  & t0   & Caller-saved temporary register \\
r8  & t1   & Caller-saved temporary register \\
r9  & t2   & Caller-saved temporary register \\
r10 & t3   & Caller-saved temporary register \\
r11 & s0   & Callee-saved register \\
r12 & s1   & Callee-saved register \\
r13 & fp   & Frame pointer / callee-saved register \\
r14 & sp   & Stack pointer \\
r15 & lr   & Link register \\
\hline
\end{tabular}
\end{center}

\section{Zero Register}

\begin{itemize}
  \item Register \(r0\) always reads as zero.
  \item Writes to \(r0\) are ignored.
\end{itemize}

\section{Return Registers}

\begin{itemize}
  \item Registers \(r1\) and \(r2\) are used for function return values.
  \item Registers \(r1\) and \(r2\) may also be used as scratch registers by callers.
  \item Registers \(r1\) and \(r2\) form the ABI return-value pair.
  \item When a function returns a 64-bit integer result, \(r1\) contains the low 32 bits and \(r2\) contains the high 32 bits.
  \item Division helper routines return the quotient in \(r1\) and the remainder in \(r2\).
\end{itemize}

\section{Argument Registers}

\begin{itemize}
  \item Registers \(r3\) through \(r6\) are used to pass the first four function arguments.
  \item These registers are named \(a0\) through \(a3\).
\end{itemize}

\section{Caller-Saved Registers}

\begin{itemize}
  \item Registers \(r1\) through \(r10\) are caller-saved.
  \item This includes:
  \begin{itemize}
    \item \(rv0\);
    \item \(rv1\);
    \item \(a0\)--\(a3\);
    \item \(t0\)--\(t3\).
  \end{itemize}
  \item A caller must save these registers before a call if it needs their values after the call returns.
\end{itemize}

\section{Callee-Saved Registers}

\begin{itemize}
  \item Registers \(r11\) through \(r13\) are callee-saved.
  \item These registers are named \(s0\), \(s1\), and \(fp\).
  \item A callee must restore these registers before returning if it modifies them.
\end{itemize}

\section{Stack Pointer}

\begin{itemize}
  \item Register \(r14\), named \(sp\), is the stack pointer.
\end{itemize}

\section{Link Register}

\begin{itemize}
  \item Register \(r15\), named \(lr\), is the link register.
  \item Call instructions write the return address to \(lr\).
  \item A function that calls another function must preserve its incoming link register value if it needs to return to its own caller.
\end{itemize}

\section{Frame Pointer}

\begin{itemize}
  \item Register \(r13\), named \(fp\), may be used as a frame pointer.
  \item Functions that do not require a frame pointer may use \(r13\) as a callee-saved general-purpose register.
\end{itemize}
